% 第5章 5.1.1：仿射核空间 H={a+bt}、<a+bt,a'+b't>=aa'+bb' 中再生关系的几何。
% (a) 系数平面 (a,b) 与 H 等距：f=1+t 对应 (1,1)；核截面 K_x=1+xt 对应 (1,x)，全部落在直线 a=1 上。
%     K_{1/2}=(1,0.5)，<f,K_{1/2}>=1.5=f(1/2)；f 在 K_{1/2} 方向上的投影长度 1.5/sqrt(1.25)=1.342。
%     K_{-1}=(1,-1) 与 f 正交，故 f(-1)=0。
% (b) 函数图像 f(t)=1+t，标出 f(1/2)=1.5 与 f(-1)=0。
\begin{tikzpicture}[font=\figurefont\footnotesize,>={Stealth[length=2.2mm]}]
\begin{scope}[x=17mm,y=17mm]
\node[anchor=south] at (0.65,1.55) {(a) 系数平面};
\draw[FigureGrid,line width=.4pt] (-0.35,-1.25) grid[step=0.5] (1.65,1.4);
\draw[->,FigureStroke,line width=.6pt] (-0.35,0) -- (1.7,0) node[right,text=FigureInk] {$a$};
\draw[->,FigureStroke,line width=.6pt] (0,-1.25) -- (0,1.45) node[above,text=FigureInk] {$b$};
% line of kernel sections
\draw[FigureGreen,line width=.9pt,dashed] (1,-1.25) -- (1,1.4) node[above,text=FigureGreen,font=\FigureNoteFont\scriptsize] {$\{K_x\}$};
% projection of f on K_{1/2}
\coordinate (P) at ({1.5/1.25},{0.75/1.25});
\draw[FigureMuted,line width=.6pt,densely dotted] (1,1) -- (P);
\draw[FigureYellow,line width=2.2pt,opacity=.8] (0,0) -- (P);
% vectors
\draw[->,FigureBlue,line width=1.1pt] (0,0) -- (1,1) node[above left=-0.5mm,text=FigureBlue] {$f$};
\draw[->,FigureGreen,line width=1.1pt] (0,0) -- (1,0.5) node[below right=-0.6mm,text=FigureGreen] {$K_{1/2}$};
\draw[->,FigureGreen,line width=1.1pt] (0,0) -- (1,-1) node[below right=-0.6mm,text=FigureGreen] {$K_{-1}$};
\draw[FigureMuted,line width=.5pt] ($(0,0)!2.2mm!(1,1)$) -- ++(0.13,-0.13) -- ($(0,0)!2.2mm!(1,-1)$);
\fill[FigureBlue] (1,1) circle (0.6pt);
\node[font=\FigureNoteFont\scriptsize,text=FigureInk,anchor=north] at (0.62,-1.27) {$\langle f,K_{-1}\rangle=0$};
\end{scope}
\begin{axis}[at={(50mm,0)},anchor=south west,yshift=-21mm,width=46mm,height=50mm,clip=false,
  axis lines=middle,xmin=-1.15,xmax=1.2,ymin=-0.3,ymax=2.2,
  xtick={-1,-0.5,0.5,1},ytick={0.5,1,1.5,2},xlabel={$t$},ylabel={$f(t)$},
  x axis line style={draw=FigureStroke,line width=.6pt,-{Stealth[length=2mm]}},y axis line style={draw=FigureStroke,line width=.6pt,-{Stealth[length=2mm]}},
  tick style={draw=FigureStroke},tick label style={font=\figurefont\scriptsize,text=FigureInk},
  label style={font=\figurefont\scriptsize},
  xlabel style={anchor=west},ylabel style={anchor=south},
  ymajorgrids,grid style={draw=FigureGrid,line width=.4pt},
  title={(b) 函数$f(t)=1+t$},title style={font=\figurefont\footnotesize,yshift=1mm}]
\addplot[draw=FigureBlue,line width=1.1pt,domain=-1:1] {1+x};
\addplot[draw=FigureYellow,line width=1.1pt] coordinates {(0.5,0) (0.5,1.5)};
\addplot[only marks,mark=*,mark size=1.5pt,FigureBlue] coordinates {(0.5,1.5) (-1,0)};
\node[anchor=west,font=\FigureNoteFont\scriptsize,text=FigureInk] at (axis cs:0.56,1.22) {$f(\tfrac12)=1.5$};
\end{axis}
\end{tikzpicture}
